\section{Discussion}\label{sec:disc}

\subsection{What the Design Buys and What It Costs}
Compared with tokens and CRLs, the central gain is the removal of an implicit-trust window: on the audited path a revoked device loses access at the next call, and on the fast path within the replica lag, a bound the operator can set and measure (Propositions~\ref{prop:revoke} and~\ref{prop:stale}; 33--147 accepted requests for the baselines against none or $\rho\Delta$). Quorum governance removes the single-registrar insider (T7) at a gas cost that scales with the voter group rather than the fleet. The tiered policy concentrates verification effort where a wrong decision hurts: at T2--T3 the gate and strict mode cut a spoofing attacker's acceptance from about 100\% to 0\% and from about half to zero, and the cost is a few percent of legitimate requests denied when the process is upset. In contrast to schemes that authenticate once and trust the session \cite{wang2026medicaliot}, nothing here is trusted by default; in contrast to conceptual smart-factory designs \cite{farhan2026ztfsf}, the policy is an executable contract with measured cost.

The costs are equally concrete. The fast path doubles per-request compute (two signature checks), the audited path is bounded by block gas (about 424 decisions per 30M-gas block by gas, and 11--18 per second measured on a single host and 31--50 per second with validators in four regions), recall of very large fleets takes tens of blocks, and the whole design adds a consortium chain, a replica and an oracle to operate. For a plant that has one owner and a good internal PKI, this is probably unjustified; the argument for it is multi-party plants in which revocation authority and audit trails must be shared.

\subsection{Deployment Guidance}
\emph{Tier assignment.} Tiers should follow consequence, not protocol: any action whose misuse could damage equipment or product, or alter safety-relevant logic, belongs on the audited, strict path; high-rate reads belong on the fast path. \emph{Gate calibration.} The anomaly gate should be set from a plant's own distribution of legitimate deviations, including start-up and batch changes (Section~\ref{sec:gate_eval}). \emph{Replica lag.} $\Delta$ should be measured, monitored and alarmed, since Proposition~\ref{prop:stale} is only as good as A5. \emph{Recall rehearsal.} Batch revocation should be rehearsed with quarantine (\texttt{suspendBatch}) first, because REVOKED is terminal and a mistaken revocation requires re-enrolment with a new key. \emph{Privacy and retention.} The ledger should hold identifiers, public keys, status, hashes and trust scores but no personal data; entries are immutable, so anything subject to erasure rules belongs off-chain and is referenced by hash. Status changes and request timing still reveal a plant's inventory and recall activity to every consortium member, so the chain should be permissioned, with read access limited to members. \emph{Keys and clocks.} Device keys belong in a secure element or TPM (A3); any attestation-freshness check needs loosely synchronised clocks in the zone (NTP or PTP) with a stated tolerance. \emph{Gas economics.} On a private PoA network gas can be zero-priced, so the figures above indicate capacity and relative cost rather than money; gasless relaying can hide fees from constrained devices \cite{madhwal2025gasless}.

\subsection{Threats to Validity and Limitations}
\emph{Internal validity.} In the main experiments devices, signals and attacks are synthetic and chosen by us, and the legitimate and attacking anomaly ranges overlap by construction. In the real-data studies only the anomaly signal is real: posture, network context and reliability are synthetic, the assignment of a flow score to a device is arbitrary, and the held-out-type test approximates but is not a real zero-day setting. The datasets have known weaknesses (duplicates, normal and attack records collected on different days, attacks from a few hosts), so the supervised AUC of 0.9999 is optimistic. None of this is a plant measurement. \emph{External validity.} The simulations ran on one two-vCPU container; no measurement was made on a microcontroller or on industrial protocols. The multi-node results come from one single-host run and one four-region run on one provider, on one day, without injected delay or partitions, and the fault tests used a single pending transaction and small samples. The replica lag of about 0.2\,s, the throughput lower bound of 50 tx/s and the failure behaviour have to be measured again on the intended deployment. The legacy-device proxy is a design with a modelled class, not an implemented gateway. \emph{Construct and conclusion validity.} Simulated exposure is counted on a logical clock and ignores propagation and finality time; the Besu runs supply the finality component but not the delay of an industrial network. Intervals come from 10--30 runs; the 500-trial attack tests have a 0.6\% rule-of-three bound.

\emph{Comparison.} No earlier blockchain-IoT or Zero Trust scheme was re-implemented, so the comparison in Table~\ref{tab:relatedwork} is qualitative, and the quantitative baselines (token, CRL, OCSP, central decision point, authenticate-once session, static-threshold policy) are our own class models. Signing and verification were not measured on a constrained device.

\emph{Trust assumptions.} The PEP is trusted for request-signature checking in Track A, and a single oracle and registrar remain single points of compromise (A4); a threshold-signed oracle \cite{hassan2023oracle} and a quorum registrar set would remove them. Multi-PEP zones require binding the PEP identity into requests (Proposition~\ref{prop:replay}). Trust scores appear in transaction data and are visible to consortium members; differential-privacy protection of exposed scores \cite{ren2026ditdp} is not implemented.

\emph{Detection.} A detector-dependent gate inherits the detector's blind spots on unseen attacks and is exposed to an adaptive attacker who can steer its own score. Adversarially robust detection, rate limiting of low-score probing and joint use of several detectors are natural defences that we did not evaluate; adversarial-AI mimicry is richer than our populations \cite{usama2026multilayer}, and the signal producers (attestation, process-aware anomaly detection, possibly cross-site learning \cite{tom2025flcti}) are inputs we assume. Learning the hand-tuned policy parameters from plant data \cite{salim2026ztblociot} is open, as is a field study joining per-device telemetry to the traffic of the same device.

\emph{Contracts.} The contracts are research code. They pass 20 tests, and a Slither run on the earlier version found no serious issue after one hardening fix (it was not repeated after the governance change), but they have not been audited or formally verified; the literature on smart-contract quality assurance \cite{wei2024quality} should be followed before any production use.
